\section{Joshua Bloch}
Professor na CMU; foi chief Java architect no Google e distinguished
engineer na Sun, onde liderou o Collections Framework e as mudanças de
linguagem do JDK~5.0 (fonte: \citeobra{bio:bloch-informit}).

\begin{itemize}
\obra{livro:bloch2017}{Noventa itens de boas práticas: imutabilidade,
  \texttt{equals}/\texttt{hashCode}, genéricos, lambdas e streams. É o livro
  de Java mais citado nesta apostila.}
\obra{livro:goetz2006}{O clássico de concorrência, com Goetz, Peierls,
  Bowbeer, Holmes e Lea. Continua valendo com virtual threads.}
\obra{livro:bloch2005}{Quebra-cabeças curtos, com Neal Gafter, que expõem
  cantos estranhos da linguagem e da JVM. Bons para abrir uma aula.}
\obra{video:bloch-api-design}{Princípios de design de API, com exemplos bons
  e ruins tirados do próprio JDK.}
\obra{artigo:bumper-sticker-api}{As máximas da palestra em uma página, como
  \enquote{fácil de usar, difícil de usar errado}.}
\obra{artigo:binary-search-broken}{O overflow em \texttt{(low + high) / 2}
  que ficou anos no JDK: uma lição sobre invariantes e testes.}
\obra{video:bloch-effective-java-3e}{Os destaques da terceira edição do
  \emph{Effective Java}, com foco em lambdas e streams.}
\end{itemize}
